\section{\textcolor{Plum}{Vollständiger Aufgabenkatalog des Usability-Tests}}
\label{Add:Usability_Test:Aufgabenkatalog}

\begin{table}[]
    \centering
    \begin{tabular}{ll}
        Aufgabe 1   & Haben wir heute einen Green Day im nicht-HA-Modus \\
                    & der Dev-Version (qe8)?\\
                    &  Wenn nein, warum?\\
        Aufgabe 2   & Wie viele BUILD-Fehler weisen die Dev-Versionen auf?\\
        Aufgabe 3.1 & Wie viele Fehler weist die [test\_cat\_3] in der älteren\\
                    & releasten Version [Nr] auf?\\
        Aufgabe 3.2 & ... und in der neueren [Nr]?\\
        Aufgabe 4.1 & Wie viele Tests (Test-Kategorien) sind in der (Dev)\\
                    & [test\_cat\_4] im [Modus 1] noch nicht gelaufen?\\
        Aufgabe 4.2 & ... und im [Modus 2]? \\
    \end{tabular}
\end{table}

% *******************************************************
\section{\textcolor{Salmon}{Begründung zur Auswahl und Reihenfolge der Aufgaben des Usability-Tests}}
\label{Add:Usability_Test:Begündung_Filterung}
\textcolor{red}{TODO: unvollständig}\\

\begin{itemize}
    \item alle erwarteten Daten können durch eine klare Zahl im neuen System abgelesen werden
    \item es muss im neuen System nicht gezählt oder zusammengerechnet werden
    \item die Farbgestaltung im neuen System ist einheitlich gewählt
    \item es wurden verschiedene nicht-einheitliche Methoden der Informationsdarstellung abgefragt
\end{itemize}

\noindent
$\Rightarrow$ Alle oben genannten Gründe sind der angenommene Idealzustand. Ob das erfolgreich umgesetzt wurde zeigen die Ergebnisse (s. \autoref{ch:Ergebnisse})\\

\textcolor{cyan}{
- Überblick über Anforderungen an das bestehende System / das neue System\\
-> wie kann man die testen?\\
- sind das Tests dabei, die nicht von den Fragen abgedeckt wurden?\\
Er umfasst den täglich benötigten allgemeinen Status der Tests, die Abfrage verschiedener Varianten von absoluten Werten, dessen Ergebnisse jeweils unterschiedlich abzulesen sind und von verschiedenen Systemen nicht zu verwechseln sind, sowie den Wechsel zwischen den beiden Dashboards. 
}

\noindent
Aufgabe 1
\begin{itemize}
    \item ein einfacher Wert zum Ablesen
    \item Start einer Diskussion, da die Aufgabe mögliche 2 Lösungen hat
\end{itemize}

\noindent
Aufgabe 2
\begin{itemize}
    \item Abfrage eines weiteren Werts
    \item Addition von klaren Zahlen im alten System
\end{itemize}

\noindent
Aufgabe 3.1
\begin{itemize}
    \item Wechsel der allgemeinen Oberfläche, um zu testen ob die neue Ansicht mit dem 2. Dashboard intuitiv verstanden wird
    \item alte Oberfläche erfordert das Zusammenzählen von 'Err' oder 'OK'\\
    -> keine Zahlenangabe, wie viele Tests gut / nicht gut gelaufen sind
\end{itemize}

%%%
den Wechsel zwischen den beiden Dashboards zum Testen der Merk- und Navigationsfähigkeit (Erlernbarkeit)
%%%

\noindent
Aufgabe 3.2
\begin{itemize}
    \item Aufmerksam machen, wie man auf das gleiche Ergebnis für das andere System gelangt
\end{itemize}

\noindent
Aufgabe 4.1
\begin{itemize}
    \item Wechsel der allgemeinen Oberfläche zurück zur ersten verwendeten Oberfläche
\end{itemize}

\noindent
Aufgabe 3.2
\begin{itemize}
    \item Aufmerksam machen, wie man auf das gleiche Ergebnis für das andere System gelangt
\end{itemize}
